\begin{question}

  \begin{questionpart}
    The simplest way to derive the Schwarz inequality goes as
    follows. First, observe

    \begin{equation*}
      \tag{1.149}
      \left(
      \bra{\alpha} + \lambda^{*}\bra{\beta} \right)
      \cdot
      \left( \ket{\alpha} + \lambda\ket{\beta} \right)
      \ge 0
    \end{equation*}

    for any complex number $\lambda$; then choose $\lambda$ in such a
    way that the preceding inequality reduces to the Schwarz
    inequality.
  \end{questionpart}

  \begin{questionpart}
    Show that the equality sign in the generalized uncertainty
    relation holds if the state in question satisfies
    \begin{equation*}
      \Delta A \ket{\alpha} = \lambda \Delta B \ket{\alpha}
    \end{equation*}
     with $\lambda$ purely imaginary.
  \end{questionpart}

  \begin{questionpart}
    Explicit calculations using the usual rules of wave mechanics
    show that the wave function for a Gaussian wave packet given by

    \begin{equation*}
      \braket{x'}{\alpha} = \left(2 \pi d^2\right)^{-1/4}
      \exp[
        \frac{i\expval{p}x'}{\hbar}
        -  \frac{\left(x' - \expval{x}\right)^2}{4d^2}
        ]
    \end{equation*}

    satisfies the minimum uncertainty relation
    
    \begin{equation*}
      \sqrt{\expval{\left(\Delta x\right)^2}}
      \sqrt{\expval{\left(\Delta p\right)^2}}
      = \frac{\hbar}{2}
    \end{equation*}

    Prove that the requirement

    \begin{equation*}
      \mel{x'}{\Delta x}{\alpha}
      = \left(\text{imaginary number}\right)
      \mel{x'}{\Delta p}{\alpha}
    \end{equation*}

    is indeed satisfied for such a Gaussian wave packet, in agreement
    with (b).
  \end{questionpart}

\end{question}

\begin{purpose}
  This book doesn't have a particularly good discussion about the
  Schwarz inequality, the uncertainty relations etc.  These are
  extremely fundamental relations in quantum mechanics and deserve a
  viewing of the man behind the curtain.  The other thing we're doing
  here is integrating wave mechanics with bra/kets.  This simplistic
  explanation belies the depth of this problem so buckle up.
\end{purpose}

\begin{answer}

  \begin{answerpart}{May the Schwarz Be With [Us All]}

    Before we get into how to \emph{prove} the that the Schwarz is
    with us, we should probably show what it \emph{is}

    \begin{equation}
      \tag{1.147}
      \braket{\alpha}{\alpha}\braket{\beta}{\beta}
      \ge \abs{\braket{\alpha}{\beta}}^2
    \end{equation}
    
    Proving the Schwarz inequality starts with the observation that
    the magnitude of an inner product of a state with its own
    hermitian conjugate is always $\ge 0$ (at least in the math we
    play with here in hello quantum).  That's essentially what
    equation 1.149 is stating:

    \begin{equation}
      \begin{split}
        0
        &\le \braket{\gamma}{\gamma}
        \\
        &= \left(
        \bra{\alpha} + \lambda^{*}\bra{\beta} \right)
        \cdot
        \left( \ket{\alpha} + \lambda\ket{\beta} \right)
        \\
      \end{split}
    \end{equation}

    If you want to grind through the $3^{\text{metric-sh**-ton}}$
    components of the complex numbers, be my guest.  I'm going to take
    it on faith that when a bra and a ket like each other very much,
    they form a very special relationship and make a real number.

    We have our target, so let's ansatz a solution and see what we
    get.  We'll start by getting rid of those complex conjugates and
    move from $\braket{\alpha}{\beta}$ to
    $\abs{\braket{\alpha}{\beta}}^2$ by choosing $\lambda$ to react
    with those:

    \begin{equation}
      \text{let } \lambda \equiv \braket{\beta}{\alpha}
    \end{equation}

    \begin{equation}
      \begin{alignedat}{2}
        %% LT_1 -> & <- RT_1  &  RT_2 -> & <- LT_2
        & &
        0
        &\le \left(
        \bra{\alpha} + \lambda^{*}\bra{\beta} \right)
        \cdot
        \left( \ket{\alpha} + \lambda\ket{\beta} \right)
        \\
        & &
        &= \braket{\alpha}{\alpha}
        + \lambda\braket{\alpha}{\beta}
        + \lambda^{*}\braket{\beta}{\alpha}
        + \abs{\lambda}^{2}\braket{\beta}{\beta}
        \\
        & &
        &= \braket{\alpha}{\alpha}
        + \braket{\beta}{\alpha}\braket{\alpha}{\beta}
        + \braket{\alpha}{\beta}\braket{\beta}{\alpha}
        + \abs{\lambda}^{2}\braket{\beta}{\beta}
        \\
        & &
        &= \braket{\alpha}{\alpha}
        + \abs{\braket{\alpha}{\beta}}^2
        + \abs{\braket{\alpha}{\beta}}^2
        + \abs{\lambda}^{2}\braket{\beta}{\beta}
        \\
        & &
        &= \braket{\alpha}{\alpha}
        + 2\abs{\braket{\alpha}{\beta}}^2
        + \abs{\braket{\alpha}{\beta}}^2\braket{\beta}{\beta}
        \\
        \implies \quad & &
        0
        &\le
        \braket{\alpha}{\alpha}\braket{\beta}{\beta}
        + 2\abs{\braket{\alpha}{\beta}}^2\braket{\beta}{\beta}
        + \abs{\braket{\alpha}{\beta}}^2\braket{\beta}{\beta}\braket{\beta}{\beta}
        \\
        \implies \quad & &
        \abs{\braket{\alpha}{\beta}}^2
        &\le
        \braket{\alpha}{\alpha}\braket{\beta}{\beta}
        + 2\abs{\braket{\alpha}{\beta}}^2\braket{\beta}{\beta}
        + \abs{\braket{\alpha}{\beta}}^2
        + \abs{\braket{\alpha}{\beta}}^2\braket{\beta}{\beta}\braket{\beta}{\beta}
        \\
      \end{alignedat}
    \end{equation}

    OK, first attempt yielded something pretty ugly, but at least we
    have the right form...plus a bunch of garbage.  Let's start by
    getting rid of the the extra $\braket{\beta}{\beta}$'s by dividing our
    $\lambda$ by $\braket{\beta}{\beta}$.

    \begin{equation}
      \text{let } \lambda \equiv \frac{\braket{\beta}{\alpha}}{\braket{\beta}{\beta}}
    \end{equation}

    \begin{equation}
      \begin{alignedat}{2}
        %% LT_1 -> & <- RT_1  &  RT_2 -> & <- LT_2
        & &
        0
        &\le
        \braket{\alpha}{\alpha}
        + \lambda\braket{\alpha}{\beta}
        + \lambda^{*}\braket{\beta}{\alpha}
        + \abs{\lambda}^{2}\braket{\beta}{\beta}
        \\
        & & &=
        \braket{\alpha}{\alpha}
        + 2\frac{\abs{\braket{\alpha}{\beta}}^2}{\braket{\beta}{\beta}}
        + \frac{\abs{\braket{\alpha}{\beta}}^2\cancel{\braket{\beta}{\beta}}}{\braket{\beta}{\beta}\cancel{\braket{\beta}{\beta}}}
        \\
        \implies \quad & & 0 &\le
        \braket{\alpha}{\alpha}\braket{\beta}{\beta}
        + 2\abs{\braket{\alpha}{\beta}}^2
        + \abs{\braket{\alpha}{\beta}}^2
        \\
        \implies \quad & & \abs{\braket{\alpha}{\beta}}^2 &\le
        \braket{\alpha}{\alpha}\braket{\beta}{\beta}
        + 4\abs{\braket{\alpha}{\beta}}^2
        \\
      \end{alignedat}
    \end{equation}

    Better, but we still have some extra terms hanging around.  What
    happens if we make our $\lambda$ negative?  It should turn that
    middle term from $+2$ to $-2$ which is a difference of $4$ and,
    fingers-crossed, should zero out the factor of $4$ in our last
    attempt.  Let's find out.

    \begin{equation}
      \text{let } \lambda \equiv -\frac{\braket{\beta}{\alpha}}{\braket{\beta}{\beta}}
    \end{equation}

    \begin{equation}
      \begin{alignedat}{2}
        %% LT_1 -> & <- RT_1  &  RT_2 -> & <- LT_2
        & &
        0
        &\le
        \braket{\alpha}{\alpha}
        + \lambda\braket{\alpha}{\beta}
        + \lambda^{*}\braket{\beta}{\alpha}
        + \abs{\lambda}^{2}\braket{\beta}{\beta}
        \\
        & & &=
        \braket{\alpha}{\alpha}
        - 2\frac{\abs{\braket{\alpha}{\beta}}^2}{\braket{\beta}{\beta}}
        + \frac{\abs{\braket{\alpha}{\beta}}^2\cancel{\braket{\beta}{\beta}}}{\braket{\beta}{\beta}\cancel{\braket{\beta}{\beta}}}
        \\
        \implies \quad & & 0 &\le
        \braket{\alpha}{\alpha}\braket{\beta}{\beta}
        - 2\abs{\braket{\alpha}{\beta}}^2
        + \abs{\braket{\alpha}{\beta}}^2
        \\
        \implies \quad & & \abs{\braket{\alpha}{\beta}}^2 &\le
        \braket{\alpha}{\alpha}\braket{\beta}{\beta}
        \qed
        \\
      \end{alignedat}
    \end{equation}

    So, if we start from the concept that magnitudes are positive, we
    can ultimately find our way to the Schwarz.  May we also find
    Yogurt.  Yogurt with raspberries.
  \end{answerpart}

  \begin{answerpart}{Setting Boundaries}
    So long as we're showing compliance, let's show what we're
    complying with, and add a few utility equations in there:

    \begin{equation}
      \tag{1.98}
      \expval{A} \equiv \mel{\alpha}{A}{\alpha}
    \end{equation}

    \begin{equation}
      \tag{1.143}
      \Delta A \equiv A - \expval{A}
    \end{equation}
    
    \begin{equation}
      \tag{1.144}
      \dispersion{A}
      = \expval{(A^2 - 2A\expval{A} + \expval{A}^2)}
      = \expval{A^2} - \expval{A}^2
    \end{equation}

    \begin{equation}
      \tag{1.146}
      \dispersion{A}\dispersion{B} \ge \frac{1}{4} \abs{\expval{[A, B]}}^2
    \end{equation}

    So, in this case, we'll want to show the equality version of that
    which is to say:

    \begin{equation}
      \inlabel{boundary}
      \dispersion{A}\dispersion{B} = \frac{1}{4} \abs{\expval{[A, B]}}^2
    \end{equation}

    whenever

    \begin{equation}
      \inlabel{complex}
      \eval{\Delta A \ket{\alpha} = iy \Delta B \ket{\alpha}}_{y \in \mathbb{R}}
    \end{equation}

    Given our handy relation, let's see what happens when we
    substitute it into the left side.  Quick note here since it
    tripped me up the first time: equation \inref{complex} only
    applies when you're acting on $\ket{\alpha}$.  So if you try to
    apply it twice to $\ket{\alpha}$ then it won't (necessarily) work
    the second time.

    \begin{equation}
      \begin{split}
        \dispersion{A}\dispersion{B}
        &=
        \mel{\alpha}{\Delta A \Delta A}{\alpha}
        \expval{(\Delta B)^2}
        \\
        &=
        (\bra{\alpha}\Delta A)(\Delta A\ket{\alpha})
        \expval{(\Delta B)^2}
        \\
        &=
        (-iy\bra{\alpha}\Delta B)(iy\Delta B\ket{\alpha})
        \expval{(\Delta B)^2}
        \\
        &=
        (-iy)(iy)\mel{\alpha}{\Delta B \Delta B}{\alpha}
        \expval{(\Delta B)^2}
        \\
        &=
        y^2\mel{\alpha}{\Delta B \Delta B}{\alpha}
        \expval{(\Delta B)^2}
        \\
        &=
        (y\expval{(\Delta B)^2})^2
        \\
      \end{split}
    \end{equation}

    Let's see what happens on the right.  We're going to make use of
    an aside from the book on page 33 just after equation 1.153 which
    doesn't have an equation number associated with it, but I'll
    provide a derivation of it anyway:

    \begin{equation}
      \inlabel{deltaNoMore}
      \begin{split}
        [\Delta A, \Delta B]
        &= (A - \expval{A})(B - \expval{B})
        - (B - \expval{B})(A - \expval{A})\\
        &= (AB - A\expval{B} - \expval{A}B + \expval{A}\expval{B})
        - (BA - B\expval{A} - \expval{B}A + \expval{B}\expval{A})\\
        &= (AB - \cancel{A\expval{B}} - \expval{A}B + \expval{A}\expval{B})
        - (BA - B\expval{A} - \cancel{\expval{B}A} + \expval{B}\expval{A})\\
        &= (AB - \cancel{A\expval{B}} - \cancel{\expval{A}B} + \expval{A}\expval{B})
        - (BA - \cancel{B\expval{A}} - \cancel{\expval{B}A} + \expval{B}\expval{A})\\
        &= (AB - \cancel{A\expval{B}} - \cancel{\expval{A}B} + \cancel{\expval{A}\expval{B}})
        - (BA - \cancel{B\expval{A}} - \cancel{\expval{B}A} + \cancel{\expval{B}\expval{A}})\\
        &= AB - BA\\
        &= [A, B]\\
      \end{split}
    \end{equation}

    That allows us to use equation \inref{complex} on the right side:


    \begin{equation}
      \begin{split}
        %% LT_1 -> & <- RT_1  &  RT_2 -> & <- LT_2
        \frac{1}{4} \abs{\expval{[A, B]}}^2
        &=
        \frac{1}{4} \abs{\expval{[\Delta A, \Delta B]}}^2
        \\
        &= 
        \frac{1}{4} \abs{\mel{\alpha}{[\Delta A, \Delta B]}{\alpha}}^2
        \\
        &= 
        \frac{1}{4}
        \abs{\mel{\alpha}{\Delta A\Delta B - \Delta B \Delta A}{\alpha}}^2
        \\
        &= 
        \frac{1}{4}
        (\mel{\alpha}{\Delta A\Delta B - \Delta B \Delta A}{\alpha})^\dagger
        (\mel{\alpha}{\Delta A\Delta B - \Delta B \Delta A}{\alpha})
        \\
        &= 
        \frac{1}{4}
        (\mel{\alpha}{\Delta B\Delta A - \Delta A \Delta B}{\alpha})
        (\mel{\alpha}{\Delta A\Delta B - \Delta B \Delta A}{\alpha})
        \\
        &= 
        \frac{1}{4}
        (\mel{\alpha}{\Delta B(iy\Delta B) - (-iy\Delta B) \Delta B}{\alpha})
        (\mel{\alpha}{(-iy\Delta B)\Delta B - \Delta B(iy\Delta B)}{\alpha})
        \\
        &= 
        \frac{1}{4}
        ((iy)\mel{\alpha}{(\Delta B)^2 + (\Delta B)^2}{\alpha})
        ((-iy)\mel{\alpha}{(\Delta B)^2 + (\Delta B)^2}{\alpha})
        \\
        &= 
        \frac{1}{4}
        (y\expval{(\Delta B)^2 + (\Delta B)^2})
        (y\expval{(\Delta B)^2 + (\Delta B)^2})
        \\
        &= 
        \frac{1}{4}
        2(y\expval{(\Delta B)^2})
        2(y\expval{(\Delta B)^2})
        \\
        &= 
        (y\expval{(\Delta B)^2})
        (y\expval{(\Delta B)^2})
        \\
        &= 
        (y\expval{(\Delta B)^2})^2\qed
        \\
      \end{split}
    \end{equation}

    So this is a truly fascinating finding.  Why does equation
    \inref{complex} ensure that we minimize our uncertainty?  Are we
    rotating about an axis going into and out of the complex plane
    with an arbitrary scale-factor $y$?  Do quaternions hold the
    answers to these questions?

    To answer the ``What's going on here?'' part of this, we're going
    to revisit the entire developmental chain and also start by
    answering a different question: ``Why is the Schwarz inequality
    even a thing?''  We'll start by recalling the inequality itself,
    but with a minor twist:

    \begin{equation}
      \begin{alignedat}{2}
        & &
        \braket{\alpha}{\alpha}\braket{\beta}{\beta}
        &\ge \abs{\braket{\alpha}{\beta}}^2\\
        \implies \quad & &
        (\bra{\psi}\Delta A)(\Delta A\ket{\psi})
        (\bra{\psi}\Delta B)(\Delta B\ket{\psi})
        &\ge
        \abs{(\bra{\psi}\Delta A)(\Delta B\ket{\psi})}^2\\
        \implies \quad & &
        \dispersion{A}
        \dispersion{B}
        &\ge
        \abs{\expval{\Delta A \Delta B}}^2
      \end{alignedat}
    \end{equation}

    Recalling a little trick

    \begin{equation}
      \begin{split}
        \frac{1}{2}[A, B] + \frac{1}{2}\{A, B\}
        &= \frac{1}{2}(AB - BA) + \frac{1}{2}(AB + BA)\\
        &= \frac{1}{2}(AB - \cancel{BA} + AB + \cancel{BA})\\
        &= \frac{1}{2}(2 AB)\\
        &= AB\\
      \end{split}
    \end{equation}

    We finalize our inequality from above

    \begin{equation}
      \begin{alignedat}{2}
        & &
        \braket{\alpha}{\alpha}\braket{\beta}{\beta}
        &\ge \abs{\braket{\alpha}{\beta}}^2\\
        \implies \quad & &
        \dispersion{A}
        \dispersion{B}
        &\ge
        \abs{\expval{\Delta A \Delta B}}^2\\
        & &
        &=
        \abs{\expval{\frac{1}{2}[\Delta A, \Delta B]
            + \frac{1}{2}\{\Delta A, \Delta B\}}}^2\\
      \end{alignedat}
    \end{equation}

    Because the cross-terms of the commutation and anti-commutation
    will cancel each other out.  It works just like a complex number
    with a real and imaginary part because, well, that's exactly what
    it is.  When the Hermitian conjugate is just the negative
    (i.e. $x^\dagger = -x$) then you're dealing with something that's
    purely imaginary.  Similarly, when the value is hermitian
    (i.e. $x^\dagger = x$), you're dealing with something purely real.
    As such, normal complex number multiplying/absoluting rules
    apply. In the end, we are left with the final form of the
    uncertainty relation

    \begin{equation}
      \inlabel{actualUncertainty}
      \dispersion{A}\dispersion{B}
      \ge
      \frac{1}{4} \abs{\expval{[\Delta A, \Delta B]}}^2
      + \frac{1}{4} \abs{\expval{\{\Delta A, \Delta B\}}}^2
    \end{equation}

    Except that \emph{isn't} how it's normally written.  Normally the
    anticommutation term is omitted with the rationale that it
    ``...can only make the inequality relation stronger.''  I mean...I
    guess that's technically true...if you throw away a positive real
    term then yeah...sure...why not...you can make it less.  If you
    want to, you can also subtract $1$!  Go nuts if you like, subtract
    $2$!  What on earth is going on here?  So...probably the most
    important take-home here is

    \begin{center}
      \textbf{\large{The book is wrong.  All the books are wrong.
          Kill all humans!}}
    \end{center}

    OK, maybe the books aren't ``wrong'' per se, but they definitely
    do a very poor job here.  Let's poke around and see why.  Using
    equation \inref{deltaNoMore}, we see that for the commutation
    term, we can drop the $\Delta$'s.  So we can rewrite equation
    \inref{actualUncertainty} with some useful labels as follows

    \begin{equation}
      \begin{alignedat}{2}
        & &
        \left(\text{Variance A}\right)
        \left(\text{Variance B}\right)
        &\ge
        \abs{\text{Covariance of A/B}}^2\\
        \implies\quad&&
        \left(\text{Variance A}\right)
        \left(\text{Variance B}\right)
        &\ge
        \frac{1}{4} \abs{\expval{[A, B]}}^2
        + \frac{1}{4} \abs{\expval{\{\Delta A, \Delta B\}}}^2
      \end{alignedat}
    \end{equation}

    In classical statistical theory, we call this finding the
    Cauchy-Schwarz
    Inequality.\incite{wikipedia:CauchySchwarzInequality} It's also
    worth noting that sometimes the word ``dispersion'' is used (even
    this solutions manual uses \textbackslash dispersion vs
    \textbackslash variance) which is \emph{technically} correct
    because ``dispersion'' is the generalized concept of
    spread-outedness whereas ``variance'' is a specific computation of
    spread-outedness.  Moving on, because of equation
    \inref{deltaNoMore}, we see that one term is a statistical measure
    and one term is something more intrinsic to the operator itself.

    As we saw earlier, the anti-commutator is the real part (much like
    $z + z* = Re(z)$) and the commutator is the imaginary part.  As
    such, there is no classical analog for the imaginary part.  So why
    do people throw away the anti-commutation?  It seems like, and I
    can't really verify this, it turns out in practice to be handy to
    do so.  This move isn't one of principle, nor is it one that will
    be readily accessible to anyone reading chapter 1 of this book.
    It's a move of expedience based upon experience we don't have from
    problems we haven't seen.

    For now, probably the right way to think of it, is that there's
    clearly something \emph{interesting} about the commutation
    part...so let's pay attention to that part.

    Another thing that should be remembered is that this uncertainty
    relation doesn't actually minimize the product of the
    dispersions...especially when you just start removing terms.  Even
    the full version we (naively) derived above with the
    anti-commutation included doesn't minimize it.  There are several
    other uncertainty relations that tighten it up, whereas the
    approach listed above loosens it; that gets into the utility of
    it.  Wikipedia has a whole page of stronger uncertainty relations
    you can look at.\incite{wikipedia:StrongerUncertaintyRelations}
 
    As you can probably tell, this mistreatment was bugging me.  I've
    learned to live with it...this is me living with it...So...in
    summary...\emph{why} does the equation \inref{complex} mean that
    the inequality works out to equality?  It's because we're ensuring
    that the \emph{real} anticommutation goes to 0 and we're left only
    with the \emph{imaginary} commutation...imagine that.
  \end{answerpart}

  \begin{answerpart}{Ride the Wave [Packet]}
    This wave packet is quite close to the one worked out in the book
    in equation 1.267

    \begin{equation}
      \tag{1.267}
      \braket{x'}{\alpha} = \left[\frac{1}{\pi^{1/4}\sqrt{d}}\right]
      \exp[
        ikx'
        -  \frac{x'^2}{2d^2}
        ]
    \end{equation}

    Let's also repeat the wave packet definition from the problem just
    so we don't have to keep scrolling to the top, also let's give it
    a number.

    \begin{equation}
      \inlabel{packet}
      \braket{x'}{\alpha} = \left(2 \pi d^2\right)^{-1/4}
      \exp[
        \frac{i\expval{p}x'}{\hbar}
        -  \frac{\left(x' - \expval{x}\right)^2}{4d^2}
        ]
    \end{equation}
    
    If we compare the two side-by-side, we get:
    
    \begin{equation}
      \begin{split}
      \braket{x'}{\alpha} = \left[\frac{1}{\pi^{1/4}\sqrt{d}}\right]
      &\exp[
        ikx'
        -  \frac{x'^2}{2d^2}
        ]\\
      \braket{x'}{\alpha} =
      \frac{1}{\sqrt[4]{2}}
      \left[
        \frac{1}{\pi^{1/4}\sqrt{d}}\right]
      &\exp[
        \frac{i\expval{p}x'}{\hbar}
        -  \frac{1}{2}\frac{\left(x' - \expval{x}\right)^2}{2d^2}
        ]\\
      \end{split}
    \end{equation}

    Next up, let's recall that the primed values are variables, so we
    can replace the unprimed $\expval{x}$ with a constant just to make
    it visually simpler and also to help reenforce that this equation
    is analogous to expressing the \ac{pdf} of a normal distribution
    about $\mu$ with a mean group momentum $\expval{p}$.

    \begin{equation}
      \braket{x'}{\alpha} =
      \frac{1}{\sqrt[4]{2}}
      \left[
        \frac{1}{\pi^{1/4}\sqrt{d}}\right]
      \exp[
        \frac{ip_0x'}{\hbar}
        -  \frac{1}{2}\frac{\left(x' - x_0\right)^2}{2d^2}
        ]
    \end{equation}

    We're now quite close to being on the garden-path from the book,
    though that is hardly a satisfying way to do things.  We can
    massage it a bit further to say that $k \equiv p_0/\hbar$
    resulting in

    \begin{equation}
      \braket{x'}{\alpha} =
      \frac{1}{\sqrt[4]{2}}
      \left[
        \frac{1}{\pi^{1/4}\sqrt{d}}\right]
      \exp[
        ikx'
        -  \frac{1}{2}\frac{\left(x' - x_0\right)^2}{2d^2}
        ]
    \end{equation}

    Sakurai has already shown the following

    \begin{equation}
      \sqrt{\expval{\left(\Delta x\right)^2}}
      \sqrt{\expval{\left(\Delta p\right)^2}}
      = \frac{\hbar}{2}
    \end{equation}

    But for the sake of practice, let's go ahead and set up that
    integral.
    
    \begin{equation}
      \begin{split}
        \dispersion{x}
        &= \expval{\left(x - \expval{x}\right)^2}\\
        &=  \expval{x^2} - \expval{x}^2\\
        &= \expval{x^2} - x_0^2\\
      \end{split}
    \end{equation}

    \begin{equation}
      \begin{split}
        \expval{x^2}
        &= \int dx' x'^2 P(x')\\
        &= \int dx' x'^2 \abs{\braket{x'}{\alpha}}^2\\
        &= \int dx' x'^2 \abs{
          \frac{1}{\sqrt[4]{2}}
          \left[
            \frac{1}{\pi^{1/4}\sqrt{d}}\right]
          \exp[
            ikx'
            - \frac{1}{2}\frac{\left(x' - x_0\right)^2}{2d^2}
          ]
        }^2\\
        &=
        \frac{1}{\sqrt{2}}
        \frac{1}{\pi^{1/2}d}
        \int dx' x'^2
        \left\{
          \exp[
            ikx'
            - \frac{1}{2}\frac{\left(x' - x_0\right)^2}{2d^2}
          ]
          \exp[
            - ikx'
            - \frac{1}{2}\frac{\left(x' - x_0\right)^2}{2d^2}
          ]
          \right\}
          \\
        &=
        \frac{1}{\sqrt{2}}
        \frac{1}{\pi^{1/2}d}
        \int dx' x'^2
          \exp[- \frac{\left(x' - x_0\right)^2}{2d^2}]
          \\
      \end{split}
    \end{equation}

    From here you can use an integral table or Wolfram Alpha or
    something if you really want to.  I am, at present, content to let
    Sakurai do the heavy lifting here.  As for $\hat{p}$, that's going
    to take a bit more doing.  We will again use equation 1.144 and we
    will again assume the group momentum $\expval{p} = p_0$, so that
    still leaves the expectation of the square ($\expval{p^2}$) to
    compute.  We know that we're looking for something of the form
    $\mel{\alpha}{p}{\alpha}$.  We can't compute that directly, but we
    \emph{can} integrate that over the entire range of $x'$ to get the
    same thing.  First, though, let's start with a handy-dandy
    identity from the book

    \begin{equation}
      \tag{1.249}
      \mel{x'}{\hat{p}}{\alpha} = -i \hbar \pdv{x'} \braket{x'}{\alpha}
    \end{equation}

    \begin{equation}
      \begin{split}
        \expval{p}
        &= p_0\\
        &= \mel{\alpha}{\hat{p}}{\alpha}\\
        &= \int dx' \bra{\alpha} \left(\dyad{x'}\right) \hat{p}\ket{\alpha}\\
        &= \int dx' \braket{\alpha}{x'} \mel{x'}{\hat{p}}{\alpha}\\
        &= \int dx' \braket{\alpha}{x'} (-i \hbar \pdv{x'})\braket{x'}{\alpha}\\
      \end{split}
    \end{equation}

    That gives us a way to set up an integral using the problem setup.
    Now let's use the same approach to compute $\expval{p^2}$.

    \begin{equation}
      \begin{split}
        \expval{p^2}
        &= \mel{\alpha}{\hat{p}^2}{\alpha}\\
        &= \int dx' \braket{\alpha}{x'} \mel{x'}{\hat{p}^2}{\alpha}\\
        &= \int dx' \braket{\alpha}{x'}
        (-i \hbar \pdv{x'})^2
        \braket{x'}{\alpha}\\
        &= \int dx'
        \frac{1}{\sqrt[4]{2}}
        \left[\frac{1}{\pi^{1/4}\sqrt{d}}\right]
        \exp[
          - ikx'
          - \frac{1}{2}\frac{\left(x' - x_0\right)^2}{2d^2}
        ]
        -\hbar^2\frac{\partial^2}{\partial x'^2}
        \frac{1}{\sqrt[4]{2}}
        \left[\frac{1}{\pi^{1/4}\sqrt{d}}\right]
        \exp[
          ikx'
          - \frac{1}{2}\frac{\left(x' - x_0\right)^2}{2d^2}
        ]\\
        &=
        \frac{1}{d\sqrt{2\pi}}
        \int dx'
        \exp[
          - ikx'
          - \frac{1}{2}\frac{\left(x' - x_0\right)^2}{2d^2}
        ]
        -\hbar^2\frac{\partial^2}{\partial x'^2}
        \exp[
          ikx'
          - \frac{1}{2}\frac{\left(x' - x_0\right)^2}{2d^2}
        ]\\
      \end{split}
    \end{equation}

    Once again, feel free to grind through that if you want.  We don't
    actually need to do any of this work...it just seems like
    something that might be nice to have hanging around.

    So, going back to our problem, we need to show that for this wave
    equation

    \begin{equation}
      \mel{x'}{\Delta x}{\alpha}
      = \left(\text{imaginary number}\right)
      \mel{x'}{\Delta p}{\alpha}
    \end{equation}

    Let's start by evaluating the right hand side

    \begin{equation}
      \begin{split}
        \mel{x'}{\Delta p}{\alpha}
        &= \mel{x'}{\hat{p} - \expval{\hat{p}}}{\alpha}\\
        &= \mel{x'}{\hat{p}}{\alpha} - \mel{x'}{p_0}{\alpha}\\
        &= \left(-i \hbar \pdv{x'}\right) \braket{x'}{\alpha}
        - p_0\braket{x'}{\alpha}\\
      \end{split}
    \end{equation}

    Since the problem specifies a $\braket{x'}{\alpha}$, we can
    substitute that in and crank out the partial derivative.

    \begin{equation}
      \begin{split}
        \mel{x'}{\Delta p}{\alpha}
        &=
        \frac{1}{(2\pi)^{1/4}\sqrt{d}}
        \left\{
        \left(-i \hbar \pdv{x'}\right)
        \exp[ikx' - \frac{1}{2}\frac{\left(x' - x_0\right)^2}{2d^2}]
        - p_0\exp[ikx' - \frac{1}{2}\frac{\left(x' - x_0\right)^2}{2d^2}]
        \right\}\\
        &=
        \frac{1}{(2\pi)^{1/4}\sqrt{d}}
        \left\{
        -i \hbar
        \left(ik - \frac{2(x'-x_0)}{4d^2} \right)
        \exp[ikx' - \frac{1}{2}\frac{\left(x' - x_0\right)^2}{2d^2}]
        - p_0\exp[ikx' - \frac{1}{2}\frac{\left(x' - x_0\right)^2}{2d^2}]
        \right\}\\
        &=
        \frac{1}{(2\pi)^{1/4}\sqrt{d}}
        \left\{
        \hbar k  + i\hbar\frac{2(x' - x_0)}{4d^2}
        - p_0
        \right\}
        \exp[ikx' - \frac{1}{2}\frac{\left(x' - x_0\right)^2}{2d^2}]
        \\
        &=
        \frac{1}{(2\pi)^{1/4}\sqrt{d}}
        \left\{
        \cancel{\hbar \frac{p_0}{\hbar}}  + i\hbar\frac{2(x' - x_0)}{4d^2}
        - \cancel{p_0}
        \right\}
        \exp[ikx' - \frac{1}{2}\frac{\left(x' - x_0\right)^2}{2d^2}]
        \\
        &=
        i\hbar\frac{2(x' - x_0)}{4d^2}
        \left\{
        \frac{1}{(2\pi)^{1/4}\sqrt{d}}
        \exp[ikx' - \frac{1}{2}\frac{\left(x' - x_0\right)^2}{2d^2}]
        \right\}
        \\
        &=
        iy(x' - x_0)\braket{x'}{\alpha}\qed
        \\
      \end{split}
    \end{equation}

    Here, we have made use of the fact that $x'$ and $x_0$ are both
    real numbers and also of an earlier observation that

    \begin{equation}
      k \equiv \frac{p_0}{\hbar}
    \end{equation}

    Sakurai wants us to show equivalence to the left-side, so let's
    take a look.

    \begin{equation}
      \begin{split}
        \mel{x'}{\Delta x}{\alpha}
        &= \mel{x'}{(\hat{x} - \expval{x})}{\alpha}\\
        &= \mel{x'}{\hat{x}}{\alpha} - \mel{x'}{\expval{x}}{\alpha}\\
        &= x'\braket{x'}{\alpha} - x_0\braket{x'}{\alpha}\\
        &= (x' - x_0)\braket{x'}{\alpha}\\
      \end{split}
    \end{equation}

    As you can easily see, our relation holds and we find that the two
    differ only by an imaginary constant.

  \end{answerpart}

\end{answer}
